\documentclass[11pt,a4paper]{article}
\usepackage[margin=1in]{geometry}
\usepackage{amsmath,amssymb,amsthm}
\usepackage{algorithm}
\usepackage{algpseudocode}
\usepackage{booktabs}
\usepackage{hyperref}

\theoremstyle{definition}
\newtheorem{definition}{\textbf{Definition}}[section]
\newtheorem{theorem}{\textbf{Theorem}}[section]
\newtheorem{lemma}[theorem]{\textbf{Lemma}}
\newtheorem{remark}{\textbf{Remark}}[section]

\title{\textbf{Process Reward Models (PRM) and Step-Level Verification}}
\author{\textbf{Team Scaibu} \\ \small Email: \texttt{jaydeep.9664920749@gmail.com} \\ \small \textbf{Architecture and Core Engine Team}}
\date{\textbf{October 2026}}

\begin{document}
\maketitle

\section{Definitions and Cognitive Mental Models}

\subsection{Formal Mathematical Definition}
\textbf{Definition (Outcome vs Process Supervision and Step-Level Credit Assignment):}
In complex multi-step reasoning, traditional Outcome Reward Models (ORMs) assign a single scalar score $R_{\text{ORM}}(x, y) \in \mathbb{R}$ to the terminal answer. Outcome supervision exhibits severe credit assignment failure: it frequently rewards false reasoning chains that reach the correct answer by sheer coincidence (e.g., compensating calculation errors, or ``two wrongs making a right''), while heavily penalizing rigorous, flawless 15-step mathematical proofs that suffer from a minor arithmetic slip on the final line. \textbf{Process Reward Models (PRMs)} provide dense, localized supervision by evaluating the validity of every individual reasoning step in the solution trajectory.

Let a generated solution be decomposed into $K$ distinct logical reasoning steps $y = (s_1, s_2, \dots, s_K)$, demarcated by newline or special step delimiters (e.g. \texttt{\textbackslash n\textbackslash n}). A PRM $r_\psi$ predicts the probability of correctness for each intermediate step $k$:
{\boldmath
\begin{equation}
p_k = P_\psi(\text{step } k \text{ is correct} \mid x, s_1, \dots, s_k) = \sigma\left( r_\psi(x, s_{\le k}) \right) \in (0, 1)
\end{equation}
}
Given ground-truth step annotations $y_k \in \{0, 1\}$, the PRM is trained via step-level binary cross-entropy:
{\boldmath
\begin{equation}
\mathcal{L}_{\text{PRM}}(\psi) = - \frac{1}{K}\sum_{k=1}^K \left[ y_k \ln p_k + (1 - y_k) \ln(1 - p_k) \right]
\end{equation}
}
To guide tree search (MCTS) or solution ranking, the full-chain quality is aggregated using the \textbf{Weakest-Link Principle} (minimum step score):
{\boldmath
\begin{equation}
S_{\text{solution}} = \min_{k \in \{1, \dots, K\}} p_k
\end{equation}
}
or under independent correctness assumptions, the path product $S_{\text{prod}} = \prod_{k=1}^K p_k$. The step of first error is identified as $k^* = \min \{k \in \{1, \dots, K\} \mid p_k < \theta_{\text{err}}\}$.

\subsection{Expanded Non-Technical Definition (Intuition for Anyone)}
\textbf{Intuitive Conceptual Definition:}
\begin{enumerate}
    \item \textbf{Show Your Work}: In elementary school math, if a student writes ``$2 + 2 = 5$, and $5 \times 0 = 0$, so the answer is $0$'', an automated multiple-choice scanner marks it 100\% correct (Outcome Supervision). But a good human teacher looks at step 1 and immediately marks it with red ink: ``Wait, $2 + 2$ is not $5$!'' (Process Supervision).
    \item \textbf{Catching Hallucinations Early}: When an AI model is writing a 20-step software migration plan or medical diagnosis, a single mistake on Step 3 ruins the entire plan. An ORM only notices the crash at Step 20. A PRM sounds the alarm at Step 3, telling the model: ``Stop right here, back up, and fix Step 3!''
    \item \textbf{The Weakest Link}: A chain of logic is only as strong as its weakest link. If a mathematical proof has 9 brilliant steps but 1 broken step, the entire proof is invalid. The minimum step score ($\min_k p_k$) perfectly captures this physical law of deductive logic.
    \item \textbf{Powering Tree Search}: PRMs allow AI systems to perform deep lookahead search (like AlphaGo) over reasoning steps, exploring multiple paths and pruning dead ends the moment a step score drops.
\end{enumerate}

\subsection{Child-Friendly Mental Model (Physical Metaphor)}
Imagine walking across a suspension bridge made of wooden planks:
\begin{itemize}
    \item \textbf{Outcome Grading}: Only checking if you made it to the other side safely. If a broken plank held your weight by pure luck, you might walk across it again tomorrow and fall through!
    \item \textbf{Process Grading (PRM)}: A bridge inspector tests every single wooden plank before you take the next step.
    \item \textbf{The Red Tag}: If Plank 4 is cracked ($p_4 < 0.5$), the inspector puts a red flag on it immediately: ``Do not step on Plank 4! Step onto the reinforced plank on the left instead.''
    \item \textbf{Result}: You cross deep chasms with 100\% guaranteed safety because every single step is verified.
\end{itemize}

\section{Core Mathematical Equations and Definitions}

\subsection{Step-Level Binary Cross-Entropy Loss}
{\boldmath
\begin{equation}
\mathcal{L}_{\text{PRM}} = - \frac{1}{K}\sum_{k=1}^K \left[ y_k \ln \sigma(r_\psi(s_{\le k})) + (1 - y_k) \ln(1 - \sigma(r_\psi(s_{\le k}))) \right]
\end{equation}
}
\begin{itemize}
    \item \textbf{Formal Definition}: Discrete binary cross-entropy loss evaluated across intermediate reasoning boundaries.
    \item \textbf{What It Means}: Trains the scalar head at each step delimiter to output high probabilities for valid logical deductions and low probabilities for invalid leaps.
    \item \textbf{Why It Is Important}: Eliminates terminal credit assignment diffusion by localizing supervision directly to the exact point of error.
\end{itemize}

\subsection{Weakest-Link Composite Solution Score}
{\boldmath
\begin{equation}
S_{\min} = \min_{1 \le k \le K} p_k
\end{equation}
}
\begin{itemize}
    \item \textbf{Formal Definition}: Infimum of intermediate step correctness probabilities along the solution path.
    \item \textbf{What It Means}: Rates a multi-step solution strictly by its most dubious deduction.
    \item \textbf{Why It Is Important}: Prevents flawed arguments with charismatic concluding paragraphs from achieving high overall ranking scores.
\end{itemize}

\subsection{First Point of Failure Detection}
{\boldmath
\begin{equation}
k^* = \inf \left\{ k \in \{1, \dots, K\} \;\middle|\; p_k < \theta_{\text{err}} \right\}
\end{equation}
}
\begin{itemize}
    \item \textbf{Formal Definition}: The earliest step index where predicted validity drops below decision threshold $\theta_{\text{err}}$.
    \item \textbf{What It Means}: Identifies the root-cause reasoning failure in broken generation chains.
    \item \textbf{Why It Is Important}: Enables targeted backtracking in tree search and targeted corrective data generation in automated repair loops.
\end{itemize}

\section{Rigorous Mathematical Analysis Checklist}

\begin{itemize}
    \item \textbf{Notation}: $K$ (number of reasoning steps), $p_k \in (0, 1)$ (step correctness probability), $y_k \in \{0, 1\}$ (ground truth step label), $S_{\min}$ (weakest link score).
    \item \textbf{Epistemic Status}: Championed by Lightman et al. (OpenAI, ``Let's Verify Step by Step'', 2023), fundamental standard for mathematical LLM alignment.
    \item \textbf{Dependency Graph}: Solution $y \to \text{Segment Steps } \{s_k\} \to \text{PRM Forward } \{p_k\} \to \text{Aggregate } S_{\min} \to \text{Beam / Tree Search Guidance}$.
    \item \textbf{Dimensions}: Reasoning chain step length $K \approx 3 - 50$, scalar probability $p_k \in [0, 1]$.
    \item \textbf{Regimes}: Mathematical reasoning, formal logic proofs, code refactoring, automated scientific discovery.
    \item \textbf{Invariants}: If all step predictions match binary ground truth perfectly, $\mathcal{L}_{\text{PRM}} \approx 0.0$ and classification accuracy is $1.0$.
    \item \textbf{Isomorphisms}: PRM is isomorphic to state-value estimation in sequential Markov Decision Processes with sparse transition errors.
\end{itemize}

\section{Theoretical Theorems and Formal Proofs}

\begin{theorem}[Weakest-Link Upper Bound on Joint Solution Correctness]
Let $E_k$ denote the event that reasoning step $k \in \{1, \dots, K\}$ is mathematically correct, with conditional probability $P(E_k \mid E_{1:k-1}) = p_k$. The joint probability that the complete chain of reasoning is valid satisfies:
{\boldmath
\begin{equation}
P\left( \bigcap_{k=1}^K E_k \right) \le \min_{1 \le k \le K} p_k
\end{equation}
}
and if errors are conditionally independent, $P(\bigcap_{k=1}^K E_k) = \prod_{k=1}^K p_k \le \min_k p_k$.
\end{theorem}

\begin{proof}
Let $A = \bigcap_{k=1}^K E_k$ be the event that every step in the chain is correct.
By standard set theory, the intersection of events is a subset of any individual constituent event:
{\boldmath
\begin{equation}
A = \bigcap_{k=1}^K E_k \subseteq E_j, \quad \forall j \in \{1, \dots, K\}
\end{equation}
}
By monotonicity of probability measures:
{\boldmath
\begin{equation}
P(A) \le P(E_j), \quad \forall j \in \{1, \dots, K\}
\end{equation}
}
Because this upper bound holds simultaneously for all step indices $j \in \{1, \dots, K\}$, it holds for the infimum:
{\boldmath
\begin{equation}
P(A) \le \min_{1 \le j \le K} P(E_j) \le \min_{1 \le j \le K} p_j = S_{\min}
\end{equation}
}
Furthermore, expanding the joint probability via the chain rule:
{\boldmath
\begin{equation}
P\left( \bigcap_{k=1}^K E_k \right) = \prod_{k=1}^K P(E_k \mid E_{1:k-1}) = \prod_{k=1}^K p_k
\end{equation}
}
Since each $p_k \in [0, 1]$, multiplying probabilities can only decrease the product: $\prod_{k=1}^K p_k \le p_j$ for all $j$, and thus $\prod_k p_k \le \min_k p_k$.
This proves that the weakest-link score $S_{\min}$ serves as a rigorous, universal upper bound on solution validity.
\end{proof}

\begin{theorem}[Strict Monotonicity of Error Localization]
Let $p_k$ represent step probabilities. If step $k^*$ is the first step where $p_{k^*} < \theta_{\text{err}}$, then for any earlier step $j < k^*$, $p_j \ge \theta_{\text{err}}$, establishing that the step index $k^*$ uniquely localizes the inception of deductive failure.
\end{theorem}

\begin{proof}
Let $k^* = \inf \{k \in \{1, \dots, K\} \mid p_k < \theta_{\text{err}}\}$.
By definition of the infimum on a finite discrete ordered set $\{1, \dots, K\}$:
If the set is non-empty, $k^*$ is the minimum element.
For any index $j$ such that $1 \le j < k^*$:
{\boldmath
\begin{equation}
j \notin \left\{ k \in \{1, \dots, K\} \;\middle|\; p_k < \theta_{\text{err}} \right\}
\end{equation}
}
Therefore, $p_j < \theta_{\text{err}}$ is false, which implies:
{\boldmath
\begin{equation}
p_j \ge \theta_{\text{err}}, \quad \forall j < k^*
\end{equation}
}
This proves that all reasoning steps strictly preceding $k^*$ satisfy the validity threshold, guaranteeing that $k^*$ is the unique, un-ambiguous point of first logical departure.
\end{proof}

\section{Foundational Architectural Questions and Epistemic Meta-Questions}

\subsection{Architectural Question 1: Active Annotation via Monte Carlo Rollouts}
\textbf{Question:} How are step-level labels $y_k$ generated automatically without expensive human mathematicians?\\
\textbf{Answer:} Using automated rollout expansion (MCTS-style): from intermediate step $k$, sample $M$ diverse completions to the end of the problem. If a high fraction of completions reach the correct verified final answer, step $k$ is labeled valid ($y_k = 1$). If completions from step $k$ universally fail to reach the correct answer, step $k$ contains a fatal flaw and is labeled invalid ($y_k = 0$).

\textbf{Meta-Question:} How does PRM step-level guidance prevent exponential tree explosion in MCTS?\\
\textbf{Answer:} Unconstrained tree search over language tokens branches exponentially. PRMs allow beam search and MCTS to evaluate states only at logical step boundaries (\texttt{\textbackslash n\textbackslash n}), pruning low-scoring branches immediately. This cuts effective branching factor from hundreds of tokens down to $3 - 5$ verified candidate deductions.

\subsection{Architectural Question 2: Minimum vs Product Score Aggregation}
\textbf{Question:} When is product aggregation $\prod_k p_k$ preferred over weakest-link $\min_k p_k$?\\
\textbf{Answer:} Weakest-link is scale-free: it treats a 5-step proof and a 25-step proof equally if both have minimum step confidence $0.92$. Product aggregation penalizes length: $(0.95)^{25} \approx 0.27$, whereas $(0.95)^5 \approx 0.77$. Weakest-link is preferred when complex problems genuinely require long derivations; product aggregation is preferred when concise, direct proofs are desired.

\textbf{Meta-Question:} Can PRMs verify natural language reasoning where steps are ambiguous?\\
\textbf{Answer:} Yes. Modern models are prompted or fine-tuned to emit explicit step demarcation tags (e.g., \texttt{Step 1: ...}, \texttt{Step 2: ...}), standardizing boundary detection and making step-level reward heads completely general across law, medical triage, and multi-turn coding.

\section{Algorithmic Implementation Specification}

\begin{algorithm}[H]
\caption{Process Reward Model (PRM) Evaluation and Step Telemetry}
\begin{algorithmic}[1]
\Require Predicted step probabilities $\mathbf{p} \in [0, 1]^K$, ground truth labels $\mathbf{y} \in \{0, 1\}^K$, threshold $\theta_{\text{err}}$, mode $agg$
\Ensure PRM loss $\mathcal{L}_{\text{PRM}}$, composite score $S_{\text{comp}}$, first error step $k^*$, accuracy $\operatorname{Acc}$, scores $\mathbf{p}$
\State Initialize $bce\_accum \gets 0.0, \quad correct \gets 0, \quad k^* \gets -1, \quad \epsilon \gets 10^{-12}$
\For{$k = 0$ \textbf{to} $K-1$}
    \State $p\_clamped \gets \max(\epsilon, \min(1.0 - \epsilon, \mathbf{p}[k]))$
    \State $label \gets \mathbf{y}[k]$
    \State $bce\_accum \gets bce\_accum - (label \cdot \ln(p\_clamped) + (1 - label) \cdot \ln(1.0 - p\_clamped))$
    \State $pred\_bin \gets 1$ \textbf{if} $p\_clamped \ge \theta_{\text{err}}$ \textbf{else} $0$
    \If{$pred\_bin == label$} \State $correct \gets correct + 1$ \EndIf
    \If{$p\_clamped < \theta_{\text{err}}$ \textbf{and} $k^* == -1$}
        \State $k^* \gets k$
    \EndIf
\EndFor
\State $\mathcal{L}_{\text{PRM}} \gets bce\_accum / K, \quad \operatorname{Acc} \gets correct / K$
\If{$agg == \text{"min"}$}
    \State $S_{\text{comp}} \gets \min_{k} \mathbf{p}[k]$
\ElsIf{$agg == \text{"product"}$}
    \State $S_{\text{comp}} \gets \prod_{k=0}^{K-1} \mathbf{p}[k]$
\Else
    \State $S_{\text{comp}} \gets \frac{1}{K}\sum_{k=0}^{K-1} \mathbf{p}[k]$
\EndIf
\State \Return $(\mathcal{L}_{\text{PRM}}, S_{\text{comp}}, k^*, \operatorname{Acc}, \mathbf{p})$
\end{algorithmic}
\end{algorithm}

\section{Big-$\Theta$ Complexity Analysis}

\begin{itemize}
    \item \textbf{Binary Cross-Entropy Loss Time}: $\Theta(K)$ scalar logarithm and clamp steps.
    \item \textbf{First Point of Error Detection}: $\Theta(K)$ worst-case linear scan.
    \item \textbf{Composite Score Aggregation}: $\Theta(K)$ for min, product, or mean reduction.
    \item \textbf{Space Complexity}: $\Theta(K)$ memory buffer for step probability arrays.
    \item \textbf{Parallel Reduction Depth}: $\mathcal{D} = \Theta(\log K)$ for path score aggregation.
\end{itemize}

\section{Single Canonical Repository Reference}

The complete deterministic scalar implementation, verification suite, and unit test benchmarks are published at:
\begin{center}
\url{https://github.com/Chief-Strategist-J/llm-obs-infra}
\end{center}

\end{document}
